\subsection{Product of exponentials model}
\label{sec:poe}

The product of exponentials (PoE) model uses the joint axes at the home
configuration $q=0$. For each joint, $\omega_i$ is its unit axis direction
and $r_i$ is its physical origin. All vectors are expressed in
\texttt{base\_link}; the endpoint is \texttt{link\_6}, as in the DH model.

\textbf{Axis directions.}
The rotation axes were identified visually in a URDF view at home.
Their positive directions follow the URDF joint convention.
The fixed tilt $\gamma=\pi/6$ gives the sine and cosine components of the
tilted axes. In particular, $\omega_2=R_x(\gamma)e_y$ and
$\omega_3=\omega_5=R_x(\gamma)e_z$, where $e_x,e_y,e_z$ are the base unit vectors.

\textbf{Joint origins.}
Starting at joint 1, each origin was found by adding the relevant robot
dimensions to the preceding origin. Each offset must first be expressed
in the base frame. With lengths in metres, this gives
\begin{equation}
 \begin{aligned}
 r_1&=\tfrac12 e_z,
 & r_2&=r_1+\tfrac45 e_x+\tfrac12 e_z,\\
 r_3&=r_2+2e_x+\tfrac54\omega_2,
 & r_4&=r_3+\tfrac32 e_x+\tfrac34\omega_3,\\
 r_5&=r_4+e_x,
 & r_6&=r_5+\tfrac32 e_x.
 \end{aligned}
 \label{eq:poe-origins}
\end{equation}
For example, the offset from joint 2 to joint 3 combines a 2 m step along
the base $x$ axis and a 1.25 m step along joint 2's tilted axis.
The resulting vectors are listed in Table~\ref{tab:poe-screws}, together
with $v_i=-\omega_i\times r_i$.

\begin{table}[!ht]
\centering
\caption{Home screw data. Each vector is written as a row; $r_i$ uses the physical
joint origin. Lengths are in metres.}
\label{tab:poe-screws}
\small
\setlength{\tabcolsep}{5pt}
\renewcommand{\arraystretch}{1.55}
\begin{tabular}{@{}c c c c@{}}
\toprule
$i$ & $\omega_i^T$ & $r_i^T$ & $v_i^T$ \\
\midrule
1 & $(0,0,1)$ & $(0,0,\tfrac12)$ & $(0,0,0)$ \\
2 & $(0,\tfrac{\sqrt3}{2},\tfrac12)$ & $(\tfrac45,0,1)$
  & $(-\tfrac{\sqrt3}{2},-\tfrac25,\tfrac{2\sqrt3}{5})$ \\
3 & $(0,-\tfrac12,\tfrac{\sqrt3}{2})$ & $(\tfrac{14}{5},\tfrac{5\sqrt3}{8},\tfrac{13}{8})$
  & $(\tfrac74,-\tfrac{7\sqrt3}{5},-\tfrac75)$ \\
4 & $(1,0,0)$ & $(\tfrac{43}{10},\tfrac{5\sqrt3-3}{8},\tfrac{13+3\sqrt3}{8})$
  & $(0,\tfrac{13+3\sqrt3}{8},\tfrac{3-5\sqrt3}{8})$ \\
5 & $(0,-\tfrac12,\tfrac{\sqrt3}{2})$
  & $(\tfrac{53}{10},\tfrac{5\sqrt3-3}{8},\tfrac{13+3\sqrt3}{8})$
  & $(\tfrac74,-\tfrac{53\sqrt3}{20},-\tfrac{53}{20})$ \\
6 & $(1,0,0)$ & $(\tfrac{34}{5},\tfrac{5\sqrt3-3}{8},\tfrac{13+3\sqrt3}{8})$
  & $(0,\tfrac{13+3\sqrt3}{8},\tfrac{3-5\sqrt3}{8})$ \\
\bottomrule
\end{tabular}
\end{table}

\textbf{Forward model.}
Each pair $(\omega_i,v_i)$ defines a screw axis and its matrix form:
\begin{equation}
 \xi_i=\begin{pmatrix}\omega_i\\v_i\end{pmatrix},\qquad
 \widehat\xi_i=\begin{pmatrix}[\omega_i]_\times&v_i\\0^T&0\end{pmatrix},
 \qquad [\omega_i]_\times p=\omega_i\times p.
 \label{eq:poe-screw}
\end{equation}
The endpoint has orientation $R_x(\gamma)$ and position $r_6$ at home.
This gives its home pose $M$ and the complete forward kinematics:
\begin{equation}
 M=\begin{pmatrix}R_x(\gamma)&r_6\\0^T&1\end{pmatrix},\qquad
 \boxed{{}^0T_E(q)=e^{\widehat\xi_1q_1}e^{\widehat\xi_2q_2}
 \cdots e^{\widehat\xi_6q_6}M.}
 \label{eq:poe-fk}
\end{equation}
The screw axes stay at their home values; the exponentials are multiplied
in joint order.
